\section{More Descriptors}

\frame{
    \begin{block}{}
       \centering
        That's nice, but ... \pause
        what else can we do?
    \end{block}
}

\frame{
    \frametitle{Transform to `Nicer' Spaces}

    \begin{block}{}
        \centering
        \only<1>{\includegraphics[height=2.5in]{pd-search/dgm-set}}%
        \only<2->{\includegraphics[height=2.5in]{pd-search/dgm-to-fcns}}%
    \end{block}
}

\frame{
    \frametitle{Some Other Spaces ...}

    \begin{columns}
        \column{.32\textwidth}
        \begin{block}{Diagram}
            \centering
            {\includegraphics[height=.85in]{stat/persistence-imagedgm}}%
        \end{block}
        \pause
        \begin{block}{Landscapes}
            \centering
            {\includegraphics[height=.85in]{stat/persistence-all-landscapes}}%
        \end{block}
        \column{.32\textwidth}
        \pause
        \begin{block}{Images}
            \centering
            {\includegraphics[height=.85in]{stat/persistence-image}}%
        \end{block}
        \pause
        \begin{block}{Silhouette}
            \centering
            {\includegraphics[height=.85in]{stat/persistence-silhouette}}%
        \end{block}

        \column{.32\textwidth}
        \pause
        \begin{block}{Betti Function}
            \centering
            {\includegraphics[height=.85in]{stat/persistence-betti}}%
        \end{block}
        \pause
        \begin{block}{Euler Char.~Fcn}
            \centering
            {\includegraphics[height=.85in]{stat/persistence-ecc}}%
        \end{block}
    \end{columns}
}

\frame{
    \frametitle{Some References ...}

    \small

    \begin{enumerate}

        \item Fr\'echet Mean: Turner, K., Mileyko, Y., Mukherjee, S., \& Harer,
            J. (2014). Fréchet means for distributions of persistence diagrams.
            Discrete \& Computational Geometry, 52(1), 44-70.

        \item Images: Adams, Emerson, Kirby, Neville, Peterson,
            Shipman, Chepushtanova, Hansen, Motta, \& Ziegelmeier. (2017). Persistence images: A
            stable vector representation of persistent homology. Journal of
            Machine Learning Research, 18.

        \item Landscapes: Bubenik (2015). Statistical topological data analysis using
            persistence landscapes. J. Mach. Learn. Res., 16(1), 77-102.

        \item Silhouettes: Chazal, Fasy, Lecci, Rinaldo, \& Wasserman, (2014,
            June). Stochastic convergence of persistence landscapes and
            silhouettes. In Proc.\ 13th Annual Symposium on
            Computational Geometry (pp. 474-483).

        \item Functional Summaries: Berry, Chen, Cisewski-Kehe, \& Fasy. (2020).
            Functional summaries of persistence diagrams. Journal of Applied and
            Computational Topology, 4(2), 211-262.

    \end{enumerate}
}

\begin{frame}
    \frametitle{Silhouettes}
    \framesubtitle{Experimental Set-Up}

    \begin{block}{}
        \begin{itemize}
            \item Training data: $2$k `glands', $500$ of each type\footnote{
                    Gland and slide simulator will be released soon; email if
                    interested! Brown, Fasy, Fox, Shenfisch, Schupbach, and Quist.}\pause
            \item Distance / DTM filtration\pause
            \item Silhouette, $p=1$\pause
            \item Test set: $400$, with $100$ of each type
            \item knn classification, with $k=3$
        \end{itemize}
    \end{block}

\end{frame}

\begin{frame}<1->[label=simulated]
    \frametitle{Silhouettes}
    \framesubtitle{Results}

    \begin{block}{Berry et al. 2020}
        \centering
        \only<1>{\includegraphics[height=2.25in]{gleason/sil-mds-input}}
        \only<2->{\includegraphics[height=2.25in]{gleason/sil-mds}}
    \end{block}
\end{frame}
